\section{Implementation correspondence and proof status}
\label{sec:conformance}
The companion \code{rules.tsv} pairs every displayed named inference rule with
an implementation path/symbol or an explicitly abstract specification case.
\code{check.py} verifies that every displayed rule has one manifest entry and
that implementation symbols/paths still exist. \code{sources.sha256} identifies
the compiler sources audited here. These checks detect missing documentation
and source drift; neither is a semantics checker or proof assistant.

\subsection{Results established by the mathematical specification}
The stated proofs establish: GYO termination and alpha-acyclic recognition;
join-tree certificate soundness; projection retention; assembly correctness;
full support after two tree reductions; conditional prefix set correctness;
exact stable index coverage; lookup representation; materialized relational
realization; pure linear fusion; preservation of finite successful
Access-to-Pipeline polls; conditional adapter adequacy; and conditional contextual
composition. The finite effect analysis additionally establishes its stated
two-clause interchange and instrumentation-consistency propositions. Each
result is quantified over the domain defined in its
statement. In particular, the graph theorems concern finite compatible-value
relations; they do not depend on pretending arbitrary Rust is pure.

The source pairing is organized by semantic ownership:
\begin{center}\small
\begin{tabular}{p{.22\linewidth}p{.70\linewidth}}\toprule
Specification & Implementation owner \\\midrule
CQ grammar/scope/FV & \code{cq-ir/src/cq.rs}, \code{cq/rust.rs} \\
RA/index judgments & \code{cq-ir/src/relational_plan.rs}, \code{index_requirements.rs} \\
D-* dependencies & \code{cq-compiler/src/cq_to_relational/dependencies.rs} \\
GYO rules & \code{cq_to_relational/hypergraph.rs} \\
P-*/Y-* rules & \code{cq_to_relational/rules.rs}, \code{JoinState} \\
C-*/A-* rules & \code{cq_to_relational/clause_rules.rs}, \code{ClauseState}, \code{AtomBindings} \\
Planning policies & \code{cq_to_relational/strategy.rs}, top-level \code{plan} \\
Base/result indexes & \code{staff_r1.rs}, \code{relational_to_indexes.rs} \\
Basic/lazy/materialized storage & \code{staff_r1.rs}, \code{indexes_to_rust/lazy.rs}, \code{indexes_to_rust.rs} \\
IP-*/R-* & \code{cq-compiler/src/passes.rs} \\
Macro contexts & \code{mini-linq/src/lib.rs} \\\bottomrule
\end{tabular}
\end{center}
Paths in this table are relative to \code{crates/}; the manifest contains exact
repository-relative paths and symbols. Statement packaging, helper
constructors and deterministic strategy ranking remain separate from the
meaning of a transition. A small function named after a rule is useful only
when its actual premises and update implement that rule.

\subsection{Obligations that are not discharged by the manuscript}
\begin{obligation}[C-Rust: typed representation adequacy]
For the claimed source domain, prove that the concrete generated bindings,
storage types, lookup keys, tuple layouts, captures, clones, and output API
satisfy Rust's type/ownership rules and the representation relation.
The current language checkers retain symbolic Rust types and do not establish
this theorem. Making successful rustc elaboration a conditional premise is
legitimate for a semantic adapter lemma; it does not prove compiler type
preservation or totality on all structurally valid inputs.
\end{obligation}
\begin{obligation}[C-Resolution: quotation and aliases]
Prove the staged parse/quote/alias construction preserves the intended
resolution of ordinary identifiers, raw identifiers, macro-introduced names,
implicit macro captures, qualified paths, and host constants. Current
mixed-site/original-span construction and lexical FV tests are evidence,
not a general hygienic substitution theorem for procedural macros.
\end{obligation}
\begin{obligation}[C-Conformance: transition implementation]
Establish a simulation from each concrete compiler state/function to the
specified state and rule, including errors and unchanged state on rejected
commands. The manifest locates this work; it does not discharge it. The
current independent checker re-reads RA operations, but its shared dependency
analysis cannot independently prove D-* soundness.
\end{obligation}
Together with C-Effects in Section~\ref{sec:cursor}, these obligations prevent
a claim of complete compiler correctness, even
within a host domain for which the abstract relational theorem is proved.
An implementation can diverge from its specification despite passing every
bounded test. The specification must be revised only when its intended
semantics is changed, not simply to excuse a failing implementation.

\subsection{Refuted or missing properties of the current accepted domain}
\label{sec:gaps}
\paragraph{E1: unrestricted effect-preserving replacement is refuted.}
Counterexample~\ref{cx:effects} is both a derivation and an executable witness.
The materialized and lazy APIs differ in build/query timing and can return
different sets with stateful clauses. To support an all-Rust semantics,
the language must state observable evaluation order, prefix multiplicity,
materialization boundaries and allowed reorderings, and the passes must
preserve those choices. Merely adding more barriers or counting free names
cannot establish that. No such all-Rust theorem is claimed here.
Until that source contract is specified, E1 does not select either route as
the correct reference or classify the other as its faulty implementation.

\paragraph{E2: local well-formedness is broader than physical lowering.}
RA permits Extend with zero fresh columns. Both physical emitters currently
label that operation as a positive \code{computed()} occurrence paired with
For over \code{once(())}; Access well-formedness requires a fully bound
positive occurrence to use If. Hence the typed source
\code{answer(x) :- R(x), let () = ();} passes CQ and RA checks but both
physical routes fail to lower it. The explicit materialized entry returns a
diagnostic. The default physical entry overclaims support in its postcondition:
in debug it panics on this input; in release it returns the diagnostic.
\code{formal_domains.rs} reproduces the domain and contract mismatch. The specification names this excluded case explicitly;
repair requires an actual valid physical construction, not weakening the
access invariant until the test passes.

\paragraph{E3: general dependency regions are not graph planning.}
D-Graph chooses an entire-query route. A query with a let, computed argument,
wildcard/repeat, or general Rust guard uses prefix planning; it does not run
GYO independently on pure regions separated by dependencies. Negation and
aggregation parse but have no implemented successful semantic lowering.
WCOJ, cardinality estimation, fractional-width optimization and projected-
output-linear bounds are not implemented algorithms in this snapshot.

\paragraph{E4: exposed lower IR does not imply arbitrary continuation support.}
Direct quotation and typed pass calls are available. Public macros have fixed
successors. Parameterized CPS routing is a possible application of the
composition theorem, not an implemented configuration feature. Any future
implementation must retain each IR's own source/target and stage contract.

Section~\ref{sec:evaluation} reports the bounded tests and counterexamples
against these separate claims. A complete end-to-end theorem for the intended
arbitrary-Rust language remains unavailable until E1's semantic contract and
the C-Rust/C-Resolution/C-Conformance/C-Effects obligations are resolved.
